\section{Conclusion}

At the national level, mandating improvements in energy efficiency could be pivotal in enabling the UK to meet its Net Zero goals by 2035. The analysis reveals that large-scale investments in photovoltaic (PV) panels, as well as advanced heating and hot water technologies, would be essential for compliance. However, by adopting a strategic approach that leverages learning-by-doing and targeted investments, the UK could substantially reduce the financial burden associated with these upgrades. The study found significant variability in the estimated costs of meeting the proposed energy standards, indicating that more refined and well-communicated strategies could improve the effectiveness and uptake of energy performance improvements beyond the current RdSAP outputs on Energy Performance Certificates (EPCs). \\


The research corroborates existing empirical evidence on the existence of a green premium for enhanced energy efficiency on property values. While improving properties to an EPC rating of C is likely to have net-positive effects on both property values and energy savings, these benefits must be weighed against the substantial costs of retrofitting. These costs vary significantly depending on the age of the property and fuel poverty in a given area as for older and rural dwellings with low energy efficiency, retrofitting costs would exceed government budgets, impeding the financial viability of these investments. Consequently, landlords of privately rented properties may find it more financially prudent to sell rather than risk exposure to unstable energy and rental markets. The study estimates that approximately 500,000 properties could be at risk of sale due to these pressures, potentially triggering an 11\% decline in housing prices within the UK market. \\

The findings of this study suggest threefold benefits of implementing a minimum EPC rating policy as this single policy could lead to sectoral emissions reductions, create green premiums for homeowners, and induce a significant correction in housing prices. While this could potentially help alleviate the current housing crisis, further research is needed to estimate the drivers behind the decisions of committing to green retrofits and the extent to which certain technologies, such as solar PVs and heat pumps, could bring combined benefits for property owners. A minimum EPC C standard for privately rented homes in England and Wales, when combined with grid-wide decarbonization strategies and tailored subsidy programs, could serve as a comprehensive 'triple-win' strategy, effectively harnessing market dynamics to achieve environmental, economic, and social benefits.
